% Generated draft: requires mathematical and source review.
\documentclass[11pt]{article}
\usepackage[margin=25mm]{geometry}
\usepackage{fontspec}
\setmainfont{Latin Modern Roman}
\usepackage{amsmath,amssymb,mathtools,unicode-math}
\setmathfont{Latin Modern Math}
\usepackage{enumitem,needspace,xurl,hyperref}
\setlength{\parindent}{0pt}
\setlength{\parskip}{5pt}
\setlength{\emergencystretch}{3em}
\newcommand{\N}{\mathbb{N}}
\newcommand{\Z}{\mathbb{Z}}
\newcommand{\Q}{\mathbb{Q}}
\newcommand{\R}{\mathbb{R}}
\newcommand{\C}{\mathbb{C}}
\newcommand{\F}{\mathbb{F}}
\newcommand{\A}{\mathbb{A}}
\newcommand{\PP}{\mathbb{P}}
\newcommand{\E}{\mathbb{E}}
\newcommand{\eps}{\varepsilon}
\newcommand{\dd}{\,\mathrm{d}}
\begin{document}
\section*{UnsolvedMath ID: 30005102}
\textbf{Source problem code:} catalogue entry
\subsection{Definitions and mathematical statement}
For a positive integer \(n\), let \(A_n\) be the group of even permutations of \(\{1,\ldots,n\}\), with permutation composition as its operation. A subset \(S\subseteq A_n\) is \emph{product-free} if
\[
S\cap SS=\varnothing,
\qquad SS=\{xy:x,y\in S\}.
\]
Equivalently, there are no \(x,y,z\in S\) satisfying \(xy=z\); the elements \(x,y,z\) are not required to be distinct. Define
\[
m(n)=\max\bigl\{|S|:S\subseteq A_n,\ S\cap SS=\varnothing\bigr\}.
\]
The problem is to determine \(m(n)\) exactly as a function of \(n\).

The supplied statement specifies no lower bound on \(n\); this draft explicitly uses the positive-integer convention, whose intended small-degree scope requires confirmation. The supplied dated literature assessment labels the problem partially solved and attributes determination of the maximum and all maximizing subsets for sufficiently large \(n\) to Keevash, Lifshitz, and Minzer. It states that the unrestricted finite-degree formulation is not fully covered. No threshold or theorem statement is supplied, and this draft does not independently verify that assessment.
\subsection{Short English statement}
Determine the largest possible cardinality of a subset of \(A_n\) containing no product of two of its own elements, allowing those two elements to coincide. The supplied curation reports a solution for sufficiently large \(n\), with the unrestricted finite-degree question not fully covered.
\subsection{Sources}
\begin{itemize}
\item Local UnsolvedMath registry, supplied record with permanent ID \(30005102\), source problem code \texttt{catalogue entry}, title \emph{Largest Product-Free Subsets of Alternating Groups}.
\item Supplied report citation: \emph{Combinatorics, Probability and Computing} (2023), Oberwolfach Reports, pp.\ 1165--1237, DOI \texttt{10.4171/owr/2022/22}. \url{https://doi.org/10.4171/owr/2022/22}. The record gives proposed year 2022. Its original excerpt identifies Peter Keevash's contribution, \emph{The largest product-free subsets of the alternating groups}, as joint work with Noam Lifshitz and Dor Minzer; it does not separately identify a proposer.
\item Supplied literature reference: Keevash--Lifshitz--Minzer, product-free alternating groups. \url{https://doi.org/10.1007/s00222-024-01273-1}.
\item Supplied primary preprint: \url{https://arxiv.org/abs/2205.15191}.
\item Literature assessment checked at \texttt{2026-08-22}: supplied dated curation, not independently verified here. It labels the problem partially solved, credits Keevash, Lifshitz, and Minzer with determining the maximum and all extremizers for sufficiently large \(n\), and says the unrestricted finite-degree formulation is not fully covered. The supplied background requests independent verification of time-sensitive status and cited claims.
\end{itemize}
\end{document}
